\documentclass[11pt]{article}
\usepackage[a4paper,margin=2cm]{geometry}
\usepackage{amsmath,amssymb}
\usepackage[hidelinks]{hyperref}
\allowdisplaybreaks
\title{Qpu.General.Primes}
\author{Tsvetan Rouschev}
\date{@uuidna/qpu v1.1.1, CC-BY-NC-ND-4.0}
\begin{document}
\maketitle
\noindent\url{https://github.com/uuidna/qpu/blob/v1.1.1/lean/Qpu/General/Primes.lean}

\paragraph{\texttt{mersenne\_127}}
the Mersenne prime 2\^{}7 − 1 (127) equals mint rays − seed.
\[ \mathrm{mint}\ \mathrm{rays} - \mathrm{seed} = 127 \]
\noindent axioms: none; proof: \url{https://github.com/uuidna/qpu/blob/v1.1.1/lean/Qpu/General/Primes.lean\#L17}

\paragraph{\texttt{mersenne\_3}}
the Mersenne prime 2\^{}2 − 1 (3) equals mint coins − seed.
\[ \mathrm{mint}\ \mathrm{coins} - \mathrm{seed} = 3 \]
\noindent axioms: none; proof: \url{https://github.com/uuidna/qpu/blob/v1.1.1/lean/Qpu/General/Primes.lean\#L11}

\paragraph{\texttt{mersenne\_31}}
the Mersenne prime 2\^{}5 − 1 (31) equals bits − seed.
\[ \mathrm{bits} - \mathrm{seed} = 31 \]
\noindent axioms: none; proof: \url{https://github.com/uuidna/qpu/blob/v1.1.1/lean/Qpu/General/Primes.lean\#L15}

\paragraph{\texttt{mersenne\_7}}
the Mersenne prime 2\^{}3 − 1 (7) equals mint n − seed, and equals rays.
\[ \mathrm{mint}\ n - \mathrm{seed} = 7 \]
\noindent axioms: none; proof: \url{https://github.com/uuidna/qpu/blob/v1.1.1/lean/Qpu/General/Primes.lean\#L13}

\paragraph{\texttt{primes\_all}}
Every exhibited prime at once: the primes the lattice carries into the clay family's number-theory wing.
\[ \mathrm{coins} = 2 \land n = 3 \land \mathrm{mint}\ \mathrm{coins} - \mathrm{seed} = 3 \land \mathrm{mint}\ n - \mathrm{seed} = 7 \land \mathrm{bits} - \mathrm{seed} = 31 \land \mathrm{mint}\ \mathrm{rays} - \mathrm{seed} = 127 \land \mathrm{rays} \cdot (\mathrm{faces} - \mathrm{seed}) = 91 \land \mathrm{rays} \cdot (\mathrm{faces} - \mathrm{seed}) = \mathrm{choose}\ \mathrm{faces}\ \mathrm{coins} \land \mathrm{faces} - n = 11 \land \mathrm{faces} - \mathrm{seed} = 13 \land \mathrm{faces} - \mathrm{seed} - (\mathrm{faces} - n) = \mathrm{coins} \land \mathrm{coins} \cdot n \cdot (n + \mathrm{coins}) \cdot \mathrm{rays} = 210 \]
\noindent axioms: none; proof: \url{https://github.com/uuidna/qpu/blob/v1.1.1/lean/Qpu/General/Primes.lean\#L32}

\paragraph{\texttt{primorial\_210}}
the primorial 2 · 3 · 5 · 7 (210) equals coins · n · (n + coins) · rays.
\[ \mathrm{coins} \cdot n \cdot (n + \mathrm{coins}) \cdot \mathrm{rays} = 210 \]
\noindent axioms: none; proof: \url{https://github.com/uuidna/qpu/blob/v1.1.1/lean/Qpu/General/Primes.lean\#L29}

\paragraph{\texttt{shor\_factored}}
the Shor modulus (91) is the product of its two prime factors: rays · (faces − seed) = 7 · 13.
\[ \mathrm{rays} \cdot (\mathrm{faces} - \mathrm{seed}) = 91 \]
\noindent axioms: none; proof: \url{https://github.com/uuidna/qpu/blob/v1.1.1/lean/Qpu/General/Primes.lean\#L19}

\paragraph{\texttt{shor\_modulus\_binomial}}
the Shor modulus equals the binomial too: rays · (faces − seed) = choose faces coins.
\[ \mathrm{rays} \cdot (\mathrm{faces} - \mathrm{seed}) = \mathrm{choose}\ \mathrm{faces}\ \mathrm{coins} \]
\noindent axioms: none; proof: \url{https://github.com/uuidna/qpu/blob/v1.1.1/lean/Qpu/General/Primes.lean\#L21}

\paragraph{\texttt{twin\_gap}}
the twin-prime gap (2) is coins: (faces − seed) − (faces − n) = coins.
\[ \mathrm{faces} - \mathrm{seed} - (\mathrm{faces} - n) = \mathrm{coins} \]
\noindent axioms: none; proof: \url{https://github.com/uuidna/qpu/blob/v1.1.1/lean/Qpu/General/Primes.lean\#L27}

\paragraph{\texttt{twin\_lower}}
the lower twin prime (11) equals faces − n.
\[ \mathrm{faces} - n = 11 \]
\noindent axioms: none; proof: \url{https://github.com/uuidna/qpu/blob/v1.1.1/lean/Qpu/General/Primes.lean\#L23}

\paragraph{\texttt{twin\_upper}}
the upper twin prime (13) equals faces − seed.
\[ \mathrm{faces} - \mathrm{seed} = 13 \]
\noindent axioms: none; proof: \url{https://github.com/uuidna/qpu/blob/v1.1.1/lean/Qpu/General/Primes.lean\#L25}
\end{document}
